\documentclass[11pt,a4paper]{article}
\usepackage{fontspec}
\setmainfont{DejaVuSerif.ttf}[Path=fonts/, BoldFont=DejaVuSerif-Bold.ttf, ItalicFont=DejaVuSerif-Italic.ttf, BoldItalicFont=DejaVuSerif-BoldItalic.ttf]
\usepackage[russian]{babel}
\usepackage{amsmath,amssymb,amsthm}
\usepackage[margin=2.3cm]{geometry}
\usepackage{microtype}
\newcommand{\E}{\mathbb E}
\newcommand{\Pp}{\mathbb P}
\newcommand{\Qq}{\mathbb Q}
\newcommand{\CE}{\operatorname{CE}}
\newcommand{\MtM}{\operatorname{MtM}}
\newcommand{\NII}{\operatorname{NII}}
\title{Стационарная модель управления процентным риском банковской книги}
\author{}
\date{}
\begin{document}
\maketitle
\vspace{-1.4cm}

\paragraph{Баланс и состояние.}
Время измеряется в годах, денежные величины --- в единицах баланса.
Банк держит кредиты объёма $L_t$ и депозиты
$D_t=\nu L_t$, где $L_0>0$ и $\nu>0$ заданы.
Все кредиты амортизируются с интенсивностью $\mu>0$.
Зададим интенсивность новой выдачи и темп чистого роста книги:
\begin{equation}
 I_t=(\mu+g_0)e^{-\chi(r_t-r_0)}L_t,\qquad
 g(r)=(\mu+g_0)e^{-\chi(r-r_0)}-\mu,\qquad
 dL_t=(I_t-\mu L_t)dt=g(r_t)L_tdt.
 \label{eq:growth}
\end{equation}
Здесь $g_0\geq0$ --- рост при опорной ставке $r_0\geq0$,
$\chi\geq0$ --- чувствительность новой выдачи к ставке.
Следовательно, $L_t=L_0\exp(\int_0^t g(r_s)ds)$ зависит от истории
ставки. При $\chi=0$ возвращается прежний случай $L_t=L_0e^{g_0t}$.
Управление $f_t\in[0,1]$ есть доля \emph{новой выдачи} под
фиксированную ставку; $1-f_t$ --- доля под плавающую ставку.
Оно не совпадает с долей fixed в действующем портфеле.

Полное абсолютное состояние есть $(L_t,X_t)$, где
$X_t=(r_t,F_t,Y_t,E_t)$ содержит короткую ставку
$r_t\geq0$, непогашенный fixed-остаток $0\leq F_t\leq L_t$,
сумму годовых fixed-купонов и договорных float-спредов $Y_t\geq0$
и бухгалтерский капитал $E_t\in\mathbb R$.
Таким образом, $L_t-F_t$ --- float-остаток, а
$C_t=E_t+D_t-L_t$ --- денежная позиция, которая может быть
отрицательной. Агрегат $Y_t$ хранит эффект ставок и спредов,
зафиксированных при прошлых выдачах; хранить каждый винтаж отдельно
не требуется.

\paragraph{Ставка и условия новой выдачи.}
Под физической мерой $\Pp$ короткая ставка следует CIR-процессу
\begin{equation}
 dr_t=\kappa_P(\bar r_P-r_t)\,dt+\sigma\sqrt{r_t}\,dW_t^{\Pp}.
 \label{eq:rate}
\end{equation}
Для ценообразования нового fixed-кредита отдельно задана мера $\Qq$
с параметрами $\kappa_Q,\bar r_Q$ и той же $\sigma$.
При $u\geq0$ положим
\[
 P_Q(r,u)=\E_r^{\Qq}
 \exp\!\left(-\int_0^u r_s\,ds\right),\qquad
 A_Q(r)=\int_0^\infty e^{-\mu u}P_Q(r,u)\,du.
\]
Если единица номинала амортизируется как $e^{-\mu u}$, то её
купонный и основной потоки равны $(k+\mu)e^{-\mu u}\,du$.
Следовательно, равенство приведённой стоимости номиналу даёт
par-купон $A_Q(r)^{-1}-\mu$.
Купон нового fixed-кредита и спред нового float-кредита задаются как
\begin{equation}
 k_F(r)=A_Q(r)^{-1}-\mu+s_F,\qquad
 s_V(q)=s_V^0-\eta q,\quad q=1-f,\quad
 0\leq\eta\leq s_V^0 .
 \label{eq:pricing}
\end{equation}
Float-кредит платит текущую базовую ставку $r_t$ плюс
зафиксированный при выдаче спред $s_V(q)$.
Параметр $\eta$ отражает снижение маржи при увеличении доли float
в новой выдаче. В текущей калибровке $s_F=s_V^0$; это равенство
\emph{надбавок}, а не полных купонов.

\paragraph{Динамика книги и поток владельцу.}
Из амортизации и новой выдачи непосредственно следуют
\begin{align}
 dF_t&=\bigl[I_tf_t-\mu F_t\bigr]dt, \label{eq:F}\\
 dY_t&=\bigl[I_t\{f_tk_F(r_t)
 +(1-f_t)s_V(1-f_t)\}-\mu Y_t\bigr]dt. \label{eq:Y}
\end{align}
При коэффициенте переноса ставки в депозитную стоимость
$\beta\in[0,1]$, депозитном спреде $s_D$, добавочном спреде
финансирования отрицательной денежной позиции $s_B$ и
операционных расходах $\omega D_t$ чистый процентный доход и прибыль равны
\begin{align}
 n_t&=Y_t+r_t(L_t-F_t)+r_tC_t
       -(\beta r_t+s_D)D_t-s_B(-C_t)^+, \label{eq:nii}\\
 \Pi_t&=n_t-\omega D_t,\qquad z^+=\max(z,0). \label{eq:profit}
\end{align}
Акционер получает \emph{знаковый} поток $b_t=\alpha\Pi_t$,
$0\leq\alpha\leq1$: при $\Pi_t<0$ он покрывает убыток.
Оставшаяся доля меняет капитал:
\begin{equation}
 dE_t=(1-\alpha)\Pi_t\,dt. \label{eq:E}
\end{equation}
Жёсткого ограничения достаточности капитала нет.
Величина $E_t/L_t$ --- модельное отношение капитала к кредитам,
не регуляторный CAR с RWA в знаменателе.

\paragraph{Автономная постановка и критерий.}
Нормируем состояние на растущую книгу:
$z_t=(r_t,a_t,y_t,e_t)=(r_t,F_t/L_t,Y_t/L_t,E_t/L_t)$.
При $c=e+\nu-1$ обозначим прибыль на единицу кредитов
\begin{equation}
 p(z)=y+r(1-a)+rc-(\beta r+s_D)\nu
       -s_B(-c)^+-\omega\nu. \label{eq:p}
\end{equation}
По правилу дифференцирования частного получаем автономную систему
\begin{equation}
 \begin{aligned}
 dr_t&=\kappa_P(\bar r_P-r_t)\,dt+\sigma\sqrt{r_t}\,dW_t^{\Pp},\\
 da_t&=(\mu+g(r_t))(f_t-a_t)\,dt,\\
 dy_t&=(\mu+g(r_t))\{f_tk_F(r_t)
               +(1-f_t)s_V(1-f_t)-y_t\}\,dt,\\
 de_t&=\{(1-\alpha)p(z_t)-g(r_t)e_t\}\,dt .
 \end{aligned} \label{eq:normalized}
\end{equation}
Искомая политика есть измеримая обратная связь $f_t=\pi(z_t)\in[0,1]$.

Для случайного продолжения $Z$ введём условный энтропийный
эквивалент
\[
 \CE_\gamma(Z\mid z)=
 -\gamma^{-1}\log\E^{\Pp}[e^{-\gamma Z}\mid z],
 \qquad \CE_0(Z\mid z)=\E^{\Pp}[Z\mid z].
\]
Это \emph{рекурсивная} оценка будущих выплат, а не однократный
энтропийный эквивалент их суммы. Коэффициент риска в денежных
единицах масштабируется как $\gamma_t=\bar\gamma/L_t$,
$\bar\gamma\geq0$. Обозначим через $V(L,X)$ оптимальную
бессрочную стоимость потока владельцу. В используемой одношаговой
постановке для шага $h>0$
\begin{equation}
 V(L,X)=\sup_{f\in[0,1]}
 \left\{B_{t,t+h}^{\mathrm{PV}}(L,X,f)
 +e^{-r_th}\CE_{\gamma_t}
   \bigl(V(L_{t+h},X_{t+h})\mid L_t=L,X_t=X,f\bigr)\right\}.
 \label{eq:bellman}
\end{equation}
Здесь $B_{t,t+h}^{\mathrm{PV}}$ --- дисконтированный внутри шага
поток $b$; после шага ставка обновляется по \eqref{eq:rate}.
Терминального вознаграждения нет. Стационарная стоимость на единицу
книги $v(z)=V(L,X)/L$ удовлетворяет
\begin{equation}
 v(z)=\sup_{f\in[0,1]}
 \left\{\frac{B_{t,t+h}^{\mathrm{PV}}}{L_t}
 +e^{-rh}G(r)
 \CE_{\bar\gamma G(r)}\!
 \bigl(v(z_{t+h})\mid z_t=z,f\bigr)\right\}.
 \label{eq:normalized_bellman}
\end{equation}
Здесь $G(r)=\exp(g(r)h)$ в используемой схеме с фиксированием
ставки внутри шага. Изменение параметра внутри $\CE$ следует из
$V(L_{t+h},X_{t+h})=L_tG(r_t)v(z_{t+h})$.
Дисконтирование идёт по рыночной короткой ставке $r_t$;
дополнительной постоянной ставки дисконтирования нет.
При $h\downarrow0$ формальное HJB-уравнение, если $v$ гладкая,
имеет вид
\begin{equation}
 0=\sup_{f\in[0,1]}
 \left\{\alpha p+\mathcal L^f v-(r-g(r))v
 -\frac{\bar\gamma}{2}\sigma^2r\,(\partial_r v)^2\right\},
 \label{eq:hjb}
\end{equation}
где $\mathcal L^f$ --- генератор процесса \eqref{eq:normalized},
включая $\tfrac12\sigma^2r\,\partial_{rr}v$.
Уравнение \eqref{eq:normalized_bellman}, а не прямое решение
\eqref{eq:hjb}, реализовано в текущем численном прототипе.

\paragraph{Диагностические метрики процентного риска.}
В состоянии $(t,X)$ фиксируем $L_t,D_t,F_t,Y_t,E_t,C_t$
на следующий год, не допускаем погашений и новых выдач и
сравниваем постоянные ставки $r$ и $r+\delta$,
$\delta=0.01$. Тогда
\begin{equation}
 \Delta\NII_{1y}
 =\delta\,[L_t-F_t+C_t-\beta D_t]\cdot1\text{ год}.
 \label{eq:dnii}
\end{equation}
Рыночная стоимость \emph{существующего} баланса при немедленной
переоценке, с депозитами по номиналу, равна
\begin{equation}
 \MtM(t,X)=C_t+L_t-F_t+
       (Y_t+\mu F_t)A_Q(r_t)-D_t
       =E_t-F_t+(Y_t+\mu F_t)A_Q(r_t).
 \label{eq:mtm}
\end{equation}
Поэтому
\begin{equation}
 \Delta\MtM=(Y_t+\mu F_t)
       [A_Q(r_t+\delta)-A_Q(r_t)].
 \label{eq:dmtm}
\end{equation}
Наконец, оптимизированная метрика
\begin{equation}
 \Delta V=V(t,r_t+\delta,F_t,Y_t,E_t)
          -V(t,r_t,F_t,Y_t,E_t) \label{eq:dv}
\end{equation}
переоценивает \emph{весь будущий} поток при дальнейшей
обратной связи. Её нельзя трактовать как разность между
управляемой и неуправляемой стратегиями: для такого вклада
надо сравнить значения двух политик при том же функционале
\eqref{eq:bellman}. Все три метрики здесь измерены в денежных
единицах, но имеют разные горизонты и экономический смысл.

\end{document}
